\chapter{Referencia de la API Lua}
\section{Modelo de ejecución}
app.lua devuelve una tabla M. M.loop() es obligatorio y se ejecuta repetidamente; M.init(), M.event(event\_type) y M.cleanup() son opcionales. Sin loop se rechaza la aplicación. Ante errores de ejecución, el firmware puede pasar a una aplicación de respaldo.

El script de sensor en sensor/*.lua devuelve una tabla sensor con sensor.get(index), que devuelve un valor numérico. El firmware expone sensor globalmente. apppath contiene la ruta de la carpeta actual de la aplicación. Cargue los módulos con require("module"); para módulos opcionales puede usar pcall(require, "module").
\section{Recursos y pantalla}
loadimg y loadfont devuelven un identificador o nil; free libera imágenes y fuentes cargadas. fill rellena toda la pantalla. update envía el fotograma y después borra la pantalla en negro. drawpolygon requiere un número par de coordenadas, al menos cuatro elementos. Llame a update una vez por fotograma y libere los recursos en cleanup.
\section{Otros módulos}
system controla retardo, retroiluminación, pull-up de V2, salida de 5 V y suspensión, y lee botones, ADC y hora. onewire busca sensores DS18B20 y lee la temperatura en Celsius. nvs utiliza el espacio NVS app; get devuelve una cadena o nil, y push un booleano. ble inicializa BLE y proporciona RPM, velocidad y marcha.
\section{display}
\begin{itemize}
\item display.loadimg(path) -> img\_handle|nil
\item display.loadfont(path) -> font\_handle|nil
\item display.free(handle)
\item display.fill(color)
\item display.update()
\item display.drawimg(img, x, y)
\item display.string(font, x, y, text, color)
\item display.fillrect(x, y, width, height, color)
\item display.drawline(x1, y1, x2, y2, color)
\item display.drawcircle(x, y, r, color)
\item display.drawfillcircle(x, y, r, color)
\item display.drawroundrect(x1, y1, x2, y2, r, color)
\item display.drawarrow(x0, y0, x1, y1, w, color)
\item display.drawfillarrow(x0, y0, x1, y1, w, color)
\item display.drawpolygon(\{x1, y1, x2, y2, ...\}, color)
\item display.RED
\item display.GREEN
\item display.BLUE
\item display.WHITE
\item display.BLACK
\item display.YELLOW
\item display.CYAN
\item display.PURPLE
\item display.GRAY
\end{itemize}
\section{system}
\begin{itemize}
\item system.delay(ms)
\item system.backlight(level)
\item system.V2\_pullup(enable)
\item system.out5V(enable)
\item system.keys() -> int
\item system.adc\_ch(channel) -> number
\item system.get\_time() -> hour, minute, second
\item system.config\_sleep(enable)
\end{itemize}
\section{onewire}
\begin{itemize}
\item onewire.scan18b20() -> count
\item onewire.read18b20(index) -> temperature\_c
\end{itemize}
\section{nvs}
\begin{itemize}
\item nvs.get(key) -> string|nil
\item nvs.push(key, value\_string) -> boolean
\end{itemize}
\section{ble}
\begin{itemize}
\item ble.init() -> true
\item ble.get\_telemetry() -> rpm, speed, gear
\end{itemize}
\section{Ejemplos y recomendaciones}
El ejemplo de voltímetro usa coordenadas dentro de 160\,\texttimes\,80 píxeles y limita la anchura de la barra. El ejemplo de sensor muestra lectura ADC y suavizado, no una conversión NTC completa: sustituya la transformación ilustrativa por la calibración real. Mantenga sensor.get rápido; guarde calibraciones con nvs y convierta las cadenas a números cuando proceda. Construya rutas a partir de apppath. Busque DS18B20 una vez y después léalos en el bucle.
\begin{verbatim}
local M = {}
local font = display.loadfont(apppath .. "/../fonts/H24XB.FNT")

local function read_voltage()
return sensor.get(1)
end

function M.loop()
local v = read_voltage()
display.fill(display.BLACK)
display.string(font, 10, 30,
    string.format("%2.1f V", v), display.WHITE)
display.fillrect(10, 60,
    math.max(0, math.min(140, math.floor((v - 9.5) * 12))),
    8, display.GREEN)
display.update()
end

function M.cleanup()
display.free(font)
end

return M
\end{verbatim}
\begin{verbatim}
sensor = {}

local filt = 0.0

function sensor.get(i)
i = i or 0
system.V2_pullup(1)

local raw = system.adc_ch(i)
filt = filt + (raw - filt) / 10.0

local value = filt * 1.0
return value
end

return sensor
\end{verbatim}
